% unit: noether.p17.title.001
\section*{17. Modules dans des domaines non commutatifs, provenant notamment d'expressions différentielles et d'expressions aux différences}
\begin{center}
Emmy Noether et Werner Schmeidler, à Göttingen.\\[0.65em]
\emph{Mathematische Zeitschrift} 8 (1920), p. 1--35
\end{center}

\begin{center}
\textbf{Table des matières.}
\end{center}
\begin{center}
\begin{minipage}{0.84\linewidth}
Introduction.\\
\S~1. Généralités formelles sur les expressions différentielles et aux différences.\\
\S~2. L'anneau de polynômes général.\\
\S~3. Modules d'un seul côté et leurs groupes quotients.\\
\S~4. Réductibilité des groupes quotients en deux sous-groupes et leurs relations avec le module.\\
\S~5. Calcul avec les unités et réductibilité en $k$ sous-groupes.\\
\S~6. Théorèmes de finitude.\\
\S~7. Un cas de décomposition univoque.\\
\S~8. Décomposition additive de groupes complètement réductibles en groupes premiers. Le théorème de l'isomorphie.\\
\S~9. Rapport entre les notions \og{} isomorphe \fg{} et \og{} de même espèce \fg{}.\\
\S~10. Existence d'une infinité de décompositions d'un groupe.\\
\S~11. Relations avec les systèmes d'équations différentielles et leurs intégrales.\\
\S~12. Exemples.
\end{minipage}
\end{center}

% unit: noether.p17.intro.001
\begin{center}\textbf{Introduction.}\end{center}

% unit: noether.p17.par.001
La théorie formelle des expressions différentielles d'une variable conduit jusqu'à une décomposition en facteurs irréductibles qui est univoque en un certain sens.\footnote{Cf. E. Landau, \emph{Ein Satz über die Zerlegung homogener linearer Differentialausdrücke in irreduzible Faktoren}, \emph{Journal für die reine und angewandte Mathematik} 124 (1902), p. 115--120, et, à sa suite, A. Loewy, \emph{Über reduzible lineare homogene Differentialgleichungen}, \emph{Mathematische Annalen} 56 (1903), p. 549--584.} Mais tandis que le théorème correspondant de l'algèbre se transporte aux polynômes de plusieurs variables, cela n'est plus possible pour les expressions différentielles partielles.\footnote{Cf. un contre-exemple de Landau publié dans H. Blumberg, \emph{Über algebraische Eigenschaften von linearen homogenen Differentialausdrücken}, dissertation inaugurale, Göttingen 1912, 52 pages, p. 51--52.} Or, pour les expressions différentielles d'une variable, à côté de la représentation en produit apparaît encore une représentation comme plus petit multiple commun, où des lois analogues valent: dans deux décompositions différentes, les constituants particuliers peuvent en effet être associés deux à deux; ils sont \og{} de même espèce \fg{}.\footnote{Cf., outre les travaux cités, A. Loewy, \emph{Über vollständig reduzible lineare homogene Differentialgleichungen}, \emph{Mathematische Annalen} 62 (1906), p. 89--117; cité comme Loewy.}

% unit: noether.p17.par.002
Cette notion de plus petit multiple commun peut maintenant être comprise comme une notion de module et être ainsi transportée à $n$ variables; on obtient ainsi, dans le présent travail, la généralisation naturelle des théorèmes mentionnés. Au-delà de cela, nous remarquons encore que, des expressions différentielles, on n'utilise que le fait suivant: on peut les considérer comme des polynômes à multiplication non commutative, pour lesquels le degré du produit est égal à la somme des degrés des facteurs. Ces conditions sont satisfaites, par exemple, aussi par les expressions aux différences; les théorèmes valent donc aussi pour celles-ci, où ils n'étaient également connus jusqu'ici que pour une variable.\footnote{Cf. G. Wallenberg--A. Guldberg, \emph{Theorie der linearen Differenzengleichungen}, Leipzig et Berlin 1911.}

% unit: noether.p17.par.003
Nous développons donc en général une théorie des modules dont les éléments sont des polynômes à multiplication non commutative, et nous ramenons pour cela l'étude des modules à celle des classes modulo ces modules.\footnote{Pour le cas des domaines commutatifs, cf. W. Schmeidler, \emph{Über Moduln und Gruppen hyperkomplexer Größen}, \emph{Mathematische Zeitschrift} 3 (1919), p. 29--42.} À la différence du cas spécial commutatif, on peut certes parler ici aussi de la somme de classes modulo un module, mais non de leur produit; on peut seulement parler du produit d'une telle classe par un polynôme. La notion de \og{} groupe quotient \fg{} (historiquement \emph{Restgruppe}: groupe des classes modulo le module) doit donc être modifiée par rapport au cas commutatif. Dans le cas spécial des polynômes d'une variable, le groupe quotient est en outre \og{} fini \fg{}, c'est-à-dire qu'il ne possède qu'un nombre fini de classes modulo le module, linéairement indépendantes. En général, un tel \og{} ordre \fg{} fini n'existe pas; les groupes sont \og{} infinis \fg{}. Pour nos méthodes, toutefois, cette différence n'entre pas en considération.

% unit: noether.p17.par.004
Comme dans le cas commutatif, à la représentation d'un module comme plus petit multiple commun de modules premiers entre eux correspond le procédé plus simple de la décomposition additive du groupe quotient. Ici l'isomorphie des groupes quotients se révèle essentielle; quand on spécialise aux expressions différentielles d'une variable, cette isomorphie devient la notion de même espèce pour les modules associés, tandis que, quand on spécialise aux modules commutatifs, elle devient l'identité. D'ailleurs, la notion de même espèce se transporte aussi à $n$ variables, et son accord avec la notion d'isomorphisme peut être démontré. Que, dans la décomposition additive du groupe quotient, il n'apparaisse toujours qu'un nombre fini de termes, cela résulte du transfert du théorème de la base de Hilbert, qui repose sur la méthode de démonstration de Gordan; la méthode de démonstration originelle de Hilbert exigerait des restrictions sur les lois de multiplication (dans la formule (6), à droite, $a$ au lieu de $a_i$, ce qui est réalisé pour les expressions différentielles, mais non pour les expressions aux différences).

% unit: noether.p17.par.005
Pour parvenir maintenant à un théorème d'unicité, il faut spécialiser les modules --- comme le fait aussi Loewy --- en modules \og{} complètement réductibles \fg{}, c'est-à-dire en ceux qui sont représentables comme plus petits multiples communs de modules premiers. Pour deux décompositions différentes, les constituants sont alors ou bien identiques deux à deux, ou bien du moins à chaque constituant de l'une des décompositions correspond un constituant isomorphe de l'autre, et réciproquement; en même temps chaque décomposition contient alors au moins deux constituants isomorphes. En outre, de l'apparition de deux constituants isomorphes dans une seule et même décomposition résulte l'existence d'une infinité de décompositions, et cela vaut même sans la restriction aux modules premiers, ce qui n'était jusqu'ici pas connu même dans le cas spécial des expressions différentielles d'une variable.

% unit: noether.p17.par.006
Enfin, pour les modules provenant d'expressions différentielles, nous abordons le rapport avec les intégrales et montrons que, pour les groupes quotients finis, l'ordre est égal au nombre maximal d'intégrales linéairement indépendantes; ainsi des fonctions arbitraires apparaissent dans l'intégrale générale exactement lorsque le groupe quotient n'est pas fini. La notion selon laquelle deux modules sont de même espèce signifie ici, comme dans le cas d'une variable, une parenté \og{} birationnelle \fg{} entre les intégrales. Quelques exemples terminent le travail.

% unit: noether.p17.par.007
La première impulsion de ce travail fut donnée, il y a plusieurs années, par une question de E. Landau à E. Noether au sujet de la généralisation du théorème du produit. E. Noether remarqua alors seulement que la question relevait de la théorie des modules dans l'anneau de polynômes non commutatif, mais elle ne pouvait pas encore être abordée. Les théorèmes de Lasker sur les modules, connus alors, n'étaient pas directement transférables à cet anneau, parce que la méthode de démonstration de Lasker repose sur le théorème selon lequel un polynôme se décompose univoquement en facteurs irréductibles.\footnote{E. Noether a remarqué récemment que les théorèmes de Lasker peuvent eux aussi être établis, sans utiliser ce théorème, par des méthodes apparentées à celles employées ici.} Seules les méthodes de Schmeidler offraient la possibilité d'un transfert, que nous avons alors réalisé en commun. Dans le détail, mentionnons que la généralisation du groupe quotient, le transfert et l'application du théorème de finitude de Hilbert, ainsi que l'existence d'une infinité de décompositions, sont dus à E. Noether; la généralisation de la notion loewyenne de réductibilité complète, la notion d'isomorphisme et son rapport avec la notion de même espèce sont dus à W. Schmeidler.

% unit: noether.p17.sec.001
\subsection*{\S~1. Généralités formelles sur les expressions différentielles et aux différences.}

% unit: noether.p17.s01.par.001
1. Pour les expressions différentielles d'une variable, on a introduit une formation symbolique du produit. Soient
\[
 a_0(x)\frac{\dd^n y}{\dd x^n}+a_1(x)\frac{\dd^{n-1}y}{\dd x^{n-1}}+\cdots+a_n(x)y=A(y),
\]
\[
 b_0(x)\frac{\dd^m y}{\dd x^m}+b_1(x)\frac{\dd^{m-1}y}{\dd x^{m-1}}+\cdots+b_m(x)y=B(y).
\]
On définit alors le produit symbolique $AB(y)$ comme l'expression différentielle
\[
 a_0(x)\frac{\dd^n B(y)}{\dd x^n}+a_1(x)\frac{\dd^{n-1}B(y)}{\dd x^{n-1}}+\cdots+a_n(x)B(y).
\]
Cette formation de produit est en réalité une multiplication de deux opérateurs; on l'écrit le plus simplement en remplaçant, comme d'ordinaire, $\frac{\dd^n}{\dd x^n}$ par $\left(\frac{\dd}{\dd x}\right)^n$. Alors $AB$ devient
\[
 \left(a_0\left(\frac{\dd}{\dd x}\right)^n+\cdots+a_n\right)
 \left(b_0\left(\frac{\dd}{\dd x}\right)^m+\cdots+b_m\right)y,
\]
où, dans l'exécution du produit, valent les règles ordinaires de différentiation des sommes et des produits. Tout tel opérateur peut maintenant être considéré comme un polynôme en la seule variable $\frac{\dd}{\dd x}=\xi$, à condition de ne pas écrire explicitement la fonction à laquelle l'opérateur est appliqué et de fixer, pour la variable $\xi$, la règle de multiplication correspondant à la différentiation d'un produit:
\[
 \frac{\dd}{\dd x}\bigl(a(x)y\bigr)=a(x)\frac{\dd}{\dd x}(y)+a'(x)y,
\]
ou, dans notre notation $(n=1,m=0)$,
\begin{equation}
 \xi\cdot a=a\cdot\xi+a'.
\tag{1}
\end{equation}

% unit: noether.p17.s01.par.002
On procède tout à fait de même pour les expressions différentielles partielles linéaires
\[
 \left(\sum a_{\alpha_1\ldots\alpha_n}(x_1,\ldots,x_n)
 \frac{\p^{\alpha_1+\cdots+\alpha_n}}
 {\p x_1^{\alpha_1}\cdots\p x_n^{\alpha_n}}\right)y
\]
d'une fonction $y$. Pour les $n$ opérateurs $\frac{\p}{\p x_1},\ldots,\frac{\p}{\p x_n}$, nous introduisons des variables $\xi_1,\ldots,\xi_n$ et obtenons ainsi le polynôme
\[
 \sum a_{\alpha_1\ldots\alpha_n}(x_1,\ldots,x_n)
 \xi_1^{\alpha_1}\cdots\xi_n^{\alpha_n}.
\]
Le produit de deux polynômes se calcule alors suivant les lois
\begin{equation}
 \xi_i a=a\xi_i+a_i \qquad (i=1,\ldots,n),
\tag{2}
\end{equation}
où $a_i$ désigne la dérivée partielle de $a$ par rapport à $x_i$; les produits des $\xi$ entre eux sont commutatifs. À l'exception de la loi commutative de la multiplication, toutes les lois de l'addition et de la multiplication valent pour les polynômes ainsi définis comme dans l'algèbre ordinaire. En particulier, pour le produit de deux tels polynômes, le degré total, ainsi que le degré en chaque variable particulière, est égal à la somme des degrés correspondants des deux facteurs.

% unit: noether.p17.s01.par.003
2. Pour les expressions aux différences (cf. Wallenberg, loc. cit.) de la forme
\[
 a_x^{(0)}y_{x+n}+a_x^{(1)}y_{x+n-1}+\cdots+a_x^{(n)}y_x=A_x,
\]
\[
 b_x^{(0)}y_{x+m}+b_x^{(1)}y_{x+m-1}+\cdots+b_x^{(m)}y_x=B_x,
\]
on a une définition correspondante du produit
\[
 (AB)_x=a_x^{(0)}B_{x+n}+a_x^{(1)}B_{x+n-1}+\cdots+a_x^{(n)}B_x,
\]
qui peut être effectuée à l'aide de la règle de multiplication
\[
 (ay)_{x+1}=a_{x+1}y_{x+1}.
\]
Si nous introduisons l'opérateur $\xi$ pour le passage de $x$ à $x+1$, alors nous obtenons $\xi\{ay\}=a^*\xi\{y\}$, où $a^*=a_{x+1}$. De nouveau on peut supprimer la fonction $y$ à laquelle l'opérateur est appliqué et l'on obtient la loi
\begin{equation}
 \xi a=a^*\xi,
\tag{3}
\end{equation}
où, avec $a$, $a^*$ ne s'annule pas non plus identiquement. De façon tout à fait analogue, pour $n$ variables,
\begin{equation}
 \xi_i a=a_i\xi_i\qquad (i=1,\ldots,n),
\tag{4}
\end{equation}
où $\xi_i$ désigne le passage de $x_i$ à $x_i+1$, et où $a_i$ désigne la fonction obtenue à partir de $a$ en remplaçant $x_i$ par $x_i+1$; elle ne s'annule donc pas identiquement si $a$ ne le fait pas. Avec ces conventions, toutes les règles d'addition et de multiplication, y compris la règle concernant le degré d'un produit, valent à nouveau, à l'exception de la loi commutative.

% unit: noether.p17.s01.par.004
Les deux cas seront subordonnés, dans le paragraphe suivant, à un anneau de polynômes défini de façon tout à fait générale avec des coefficients arbitraires.

% unit: noether.p17.sec.002
\subsection*{\S~2. L'anneau de polynômes général.}

% unit: noether.p17.s02.par.001
Nous partons d'un corps $P$ défini abstraitement de manière arbitraire (dont nous désignons les éléments par des minuscules latines), pour lequel valent toutes les règles de l'algèbre, en particulier aussi la loi commutative de la multiplication avec inversion univoque. À ce corps nous adjoignons $n$ indéterminées $\xi_1,\ldots,\xi_n$; c'est-à-dire que nous considérons le domaine d'intégrité $J$ constitué par les produits $a\xi_1,b\xi_2,\ldots,c\xi_n$ et par leurs sommes finies, toutes les lois d'addition et de multiplication devant valoir à l'exception de la loi commutative de la multiplication. En particulier, nous appelons \og{} polynômes \fg{} les sommes finies de la forme
\[
 F(\xi_1,\ldots,\xi_n)=\sum a_{\nu_1\ldots\nu_\rho}\xi_{\nu_1}\cdots\xi_{\nu_\rho}
\]
(elles seront toujours désignées par des majuscules latines). À la place de la loi commutative, nous soumettons maintenant les éléments de $J$ à des lois de produit telles que \emph{le produit de deux polynômes soit de nouveau un polynôme}; le domaine de tous les polynômes de $J$ épuise donc déjà $J$. Pour cela il est manifestement nécessaire et suffisant que, pour les éléments $\xi_1,\ldots,\xi_n$ et un élément arbitraire $a$ de $P$, valent des règles de produit de la forme
\begin{equation}
 \xi_i a=F_i(\xi_1,\ldots,\xi_n)\qquad (i=1,\ldots,n),
\tag{5}
\end{equation}
où $F_i(\xi_1,\ldots,\xi_n)$ désigne un polynôme quelconque dépendant naturellement de $\xi_i$ et de $a$. (En particulier, les règles (2) et (4) pour les expressions différentielles et pour les expressions aux différences sont respectivement de cette forme.) Car, puisque $\xi_i a$ n'est pas en soi un polynôme, mais appartient par hypothèse au domaine, l'équation est nécessaire. Réciproquement, par application finiment répétée de la formule (5), chaque élément du domaine $J$ peut être transformé en un polynôme.

% unit: noether.p17.s02.par.002
Pour l'unité du corps, que nous désignons par 1, nous exigerons en outre $\xi_i\cdot1=1\cdot\xi_i=\xi_i$. Nous stipulerons aussi que la multiplication de $\xi_1,\ldots,\xi_n$ entre eux est commutative.\footnote{Cela n'est nécessaire qu'à partir de \S~6 pour le théorème de finitude et ses conséquences.} Chaque polynôme de $J$ s'écrit alors exactement comme un polynôme ordinaire en $\xi_1,\ldots,\xi_n$ à coefficients dans $P$, et deux polynômes ne coïncident que lorsque tous leurs coefficients coïncident.

% unit: noether.p17.s02.par.003
Plus précisément, dans ce qui suit,\footnote{Ceci aussi n'est utilisé qu'à partir de \S~6.} nous exigeons qu'à la place de (5) vaille la règle
\begin{equation}
 \xi_i a=a_i\xi_i+b_i\qquad (a_i\ne0)\quad(i=1,\ldots,n).
\tag{6}
\end{equation}
Ici encore les lois (2) et (4) sont en particulier incluses. Alors le degré d'un produit est égal à la somme des degrés des facteurs; ceci vaut tant pour le degré total que pour le degré en chaque variable particulière, de sorte que le produit de deux polynômes qui ne s'annulent pas identiquement est également différent de 0.

% unit: noether.p17.sec.003
\subsection*{\S~3. Modules d'un seul côté et leurs groupes quotients.}

% unit: noether.p17.s03.par.001
Pour notre domaine $J$, puisque la multiplication n'est pas commutative, les espèces suivantes de modules\footnote{Nous désignons les modules par des majuscules allemandes.} peuvent être distinguées:\footnote{Les idéaux d'un seul côté ont déjà été considérés par A. Hurwitz; cf. \emph{Vorlesungen über die Zahlentheorie der Quaternionen} (Berlin, Springer 1919), sixième leçon. Toutefois, il s'agit là essentiellement d'idéaux principaux; il en est de même des travaux de Du Pasquier cités dans la note 1 de ce livre.}

% unit: noether.p17.s03.par.002
1. Le \emph{module du côté droit}, qui, avec $M$, contient aussi $FM$ pour un polynôme arbitraire $F$, et qui, avec $M_1$ et $M_2$, contient aussi $M_1+M_2$.

% unit: noether.p17.s03.par.003
2. Le \emph{module du côté gauche}, qui, avec $M$, contient aussi $MF$, et qui, avec $M_1$ et $M_2$, contient aussi $M_1+M_2$.

% unit: noether.p17.s03.par.004
Nous désignons ces deux espèces par le nom commun de \emph{modules d'un seul côté}; les modules qui sont à la fois du côté droit et du côté gauche seront appelés \emph{bilatères}. Dans ce qui suit nous nous limitons aux modules du côté droit et remarquons que la théorie de ceux du côté gauche peut être édifiée exactement de la même façon.

% unit: noether.p17.s03.par.005
Nous disons que deux polynômes $F$ et $G$ sont congrus modulo $\mM$, et écrivons $F\equiv G(\mM)$, si leur différence est divisible par $\mM$, c'est-à-dire si elle appartient à $\mM$. L'ensemble de tous les polynômes congrus à un polynôme donné forme une \emph{classe modulo $\mM$}. (Les classes modulo $\mM$ seront désignées par de petites lettres allemandes.) Puisque de
\[
 F_1\equiv F_2(\mM)\qquad\text{et}\qquad G_1\equiv G_2(\mM)
\]
on déduit que $F_1+G_1\equiv F_2+G_2(\mM)$, la somme de deux classes modulo $\mM$, $\mx$ et $\my$, est définie et est de nouveau une classe modulo $\mM$, $\mx+\my$. En revanche, parce que le module est d'un seul côté, on ne peut pas parler du produit de deux telles classes. En effet, de
\[
 F_1=F_2+M_1,\qquad G_1=G_2+M_2
\]
on obtient seulement
\[
 F_1G_1=F_2G_2+F_2M_2+M_1G_2+M_1M_2;
\]
mais comme $M_1G_2$ n'a pas à appartenir à $\mM$ lorsque $M_1$ appartient à $\mM$, $F_1G_1$ et $F_2G_2$ ne sont pas nécessairement congrus en général. En revanche, si $M_1=0$, donc $F_1=F_2=F$, on voit que $G_1\equiv G_2$ implique $FG_1\equiv FG_2$. Ainsi une classe arbitraire $\my$ modulo $\mM$ peut être multipliée par devant par un polynôme $F$, et cela donne une classe modulo $\mM$ bien définie, $F\my=\mz$.

% unit: noether.p17.s03.par.006
Pour le calcul avec les classes modulo $\mM$ valent les lois commutative et associative de l'addition, ainsi que la loi de réversibilité univoque de l'addition; en outre, pour la multiplication de ces classes par des polynômes et l'addition, vaut la loi distributive $F(\my_1+\my_2)=F\my_1+F\my_2$, comme il résulte directement du calcul avec les polynômes modulo $\mM$.

% unit: noether.p17.s03.par.007
La classe modulo $\mM$ à laquelle appartient l'unité sera appelée classe unité, ou unité, et notée $\me$. Alors $F\me=\mx$ est la classe modulo $\mM$ à laquelle appartient le polynôme $F$. Contrairement aux classes modulo $\mM$ en général, le produit $\mx\me$ est cependant ici aussi bien défini et égal à $\mx$; car de
\[
 F_1=F_2+M_1,\qquad G_1=1+M_2
\]
et de $M_1\cdot1=M_1$ on déduit que
\[
 F_1G_1=F_2+M.
\]
L'ensemble des classes modulo $\mM$ s'appelle le \emph{groupe quotient de} $\mM$. Un groupe quotient possède donc les propriétés suivantes: 1. avec chaque classe $\mx$ et un polynôme arbitraire $F$, il contient la classe $F\mx$; et avec les classes $\mx_1$ et $\mx_2$, il contient la classe $\mx_1+\mx_2$.\footnote{Cette propriété pourrait, en un sens plus général, être appelée la \og{} propriété de module \fg{} du groupe quotient.} 2. Il possède une unité $\me$, de sorte que $F\me$ représente, pour chaque polynôme $F$, la classe $\mx$ de $F$ modulo $\mM$. \emph{Ainsi, lorsque} $F$ \emph{parcourt tous les polynômes,} $F\me$ \emph{parcourt tout le groupe quotient.}

% unit: noether.p17.sec.004
\subsection*{\S~4. Réductibilité du groupe quotient en deux sous-groupes et sa relation avec le module.}

% unit: noether.p17.s04.par.001
Pour les modules d'un seul côté, les notions de \emph{divisibilité}, de \emph{plus grand commun diviseur} et de \emph{plus petit multiple commun} demeurent les mêmes que dans le cas bilatère, comme le montrent les définitions suivantes.

% unit: noether.p17.s04.par.002
Un module $\mM$ est dit divisible par un autre module $\mN$ si $M\equiv0(\mM)$ entraîne aussi $M\equiv0(\mN)$.

% unit: noether.p17.s04.par.003
Pour deux modules $\mM$ et $\mN$, le \emph{plus grand commun diviseur} $(\mM,\mN)$ est défini comme l'ensemble de tous les polynômes représentables sous la forme $M+N$, où $M\equiv0(\mM)$ et $N\equiv0(\mN)$; c'est de nouveau un module. Il en va de même pour plusieurs modules.

% unit: noether.p17.s04.par.004
L'ensemble de tous les polynômes divisibles à la fois par $\mM$ et par $\mN$ forme un module $[\mM,\mN]$, le \emph{plus petit multiple commun} de $\mM$ et de $\mN$. Cette définition s'applique aussi de façon correspondante à plusieurs modules, et même à un ensemble de modules de cardinalité arbitraire.

% unit: noether.p17.s04.par.005
Deux modules sont dits premiers entre eux lorsque leur plus grand commun diviseur est le module unité, c'est-à-dire l'ensemble de tous les polynômes.

% unit: noether.p17.s04.par.006
Nous partons maintenant d'un module $\mM$ qui est représentable comme plus petit multiple commun de deux modules premiers entre eux $\mN$ et $\mL$, supposés tous deux être des diviseurs propres:
\begin{equation}
 \mM=[\mN,\mL].
\tag{7}
\end{equation}
Nous examinons le groupe quotient $\mG$ de $\mM$. Par hypothèse,
\begin{equation}
 1\equiv N+L(\mM)
\tag{8}
\end{equation}
pour deux polynômes $N\equiv0(\mN)$ et $L\equiv0(\mL)$. De là nous montrons
\begin{equation}
 \mL=(\mM,L),\qquad \mN=(\mM,N),
\tag{9}
\end{equation}
où, par exemple, $(\mM,L)$ désigne le module formé de tous les polynômes de la forme $M+FL$. En effet, d'abord $(\mM,L)\equiv0(\mL)$. Réciproquement, il résulte de (8), pour un polynôme arbitraire $F$, que
\begin{equation}
 F\equiv FN(\mL).
\tag{10}
\end{equation}
Si maintenant $F\equiv0(\mL)$, alors $FN\equiv0(\mL)$, donc $\equiv0(\mM)$, et par conséquent $F\equiv FL(\mM)$. On démontre de la même manière $\mN=(\mM,N)$. Il en résulte incidemment, puisque $\mN$ et $\mL$ sont des diviseurs propres de $\mM$, que ni $L\equiv0$ ni $N\equiv0(\mM)$ ne peut avoir lieu.

% unit: noether.p17.s04.par.007
Soit $\ma$ la classe de $N$ modulo $\mM$, et $\mb$ la classe de $L$ modulo $\mM$. Alors, d'après (8),
\[
 \me=\ma+\mb\qquad(\ma\ne0,\ \mb\ne0).
\]
Comme chaque classe modulo $\mM$ est de la forme $F\me$, la loi distributive donne, pour chaque classe modulo $\mM$,
\begin{equation}
 F\me=F\ma+F\mb\qquad(\ma\ne0,\ \mb\ne0).
\tag{11}
\end{equation}
Cette représentation d'une classe arbitraire modulo $\mM$ comme somme de deux classes de la forme $G\ma$ et $H\mb$ est cependant \emph{unique}. Car, à partir d'une relation $0=G\ma+H\mb$, si $K$ est un polynôme représentant la classe $G\ma=-H\mb$, alors manifestement $K\equiv GN(\mM)$, $K\equiv-HL(\mM)$, donc $K\equiv0(\mN)$, $K\equiv0(\mL)$, et par suite $K\equiv0(\mM)$; c'est-à-dire $G\ma=0$, $H\mb=0$.

% unit: noether.p17.s04.par.008
Réciproquement, supposons maintenant que la relation (11) soit satisfaite pour chaque polynôme $F$, et cela de manière unique au sens précédent. Soient alors $N$ un polynôme de $\ma$ et $L$ un polynôme de $\mb$. Nous formons les deux modules (9) et affirmons (7) et (8). Cette dernière relation résulte de (11) pour $F=1$. En outre, par la définition de $\mN$ et de $\mL$, $\mM$ est un multiple commun de ces deux modules, et même le plus petit. Car $F\equiv0(\mL)$ et $F\equiv0(\mN)$ entraînent $F\equiv HL\equiv GN(\mM)$; donc $G\ma=H\mb$, $G\ma-H\mb=0$, et par suite, par l'unicité de la représentation, $G\ma=0$, $H\mb=0$, $F\equiv0(\mM)$. Il est donc montré que la \emph{représentation d'un module} $\mM$ \emph{comme plus petit multiple commun de deux modules premiers entre eux est équivalente à la décomposition additive unique du groupe quotient} $\mG$ \emph{de} $\mM$, \emph{telle qu'elle est donnée par la formule} (11).

% unit: noether.p17.s04.par.009
Nous examinons maintenant le système $\mA$ des classes modulo $\mM$ de la forme $F\ma$, que nous voulons mettre en relation avec le groupe quotient $\mbarA$ de $\mL$. D'après (11), chaque polynôme $F$ de $\mL$ appartient à la classe $F\ma$ modulo $\mM$, tandis que $\bara$ est la classe de $N$ modulo $\mL$, et donc la classe unité de $\mbarA$. Si nous associons maintenant aux classes $F\ma$ de $\mA$ les classes $F\bara$ de $\mbarA$, alors à chaque classe de $\mA$ est assignée une classe de $\mbarA$, de telle sorte que la somme de deux classes de $\mA$ corresponde à la somme des classes assignées de $\mbarA$, et que le produit d'une classe de $\mA$ par un polynôme corresponde au produit de la classe assignée de $\mbarA$ par le même polynôme. Cette assignation est en outre biunivoque: $F\ma=0$ entraîne $F\equiv0(\mL)$ d'après (11), donc $F\bara=0$; réciproquement, $F\bara=0$, donc $F\equiv0(\mL)$, signifie aussi, d'après (10), que $FN\equiv0(\mL)$, donc $FN\equiv0(\mM)$, c'est-à-dire $F\ma=0$.

% unit: noether.p17.s04.par.010
Cela nous conduit à une notion fondamentale, que nous définissons comme suit:

% unit: noether.p17.def.001
\emph{Une correspondance biunivoque entre deux systèmes $\mA$ et $\mbarA$ de classes modulo leurs modules respectifs s'appelle isomorphe si, à la somme de deux classes de $\mA$, est assignée la somme des classes correspondantes de $\mbarA$, et si, au produit d'une classe arbitraire de $\mA$ par un polynôme, est assigné le produit de la classe correspondante de $\mbarA$ par le même polynôme.}

% unit: noether.p17.s04.par.011
Un sous-système $\mA$ de classes modulo $\mM$ dans $\mG$ qui est isomorphe au groupe quotient d'un module s'appelle un \emph{sous-groupe} de $\mG$. Comme le montre la preuve précédente, l'assignation est ici plus précisément telle que chaque classe de $\mA$ provient de la classe correspondante de $\mbarA$ par adjonction de certains polynômes, à savoir de tous les polynômes appartenant à $\mL$ mais non à $\mM$.

% unit: noether.p17.s04.par.012
De même, on montre que l'ensemble de toutes les classes modulo $\mM$ de la forme $F\mb$ constitue un sous-groupe $\mB$ de $\mG$ qui est isomorphe au groupe quotient $\mbarB$ de $\mN$.

% unit: noether.p17.s04.par.013
Un groupe quotient $\mG$ dont les classes admettent une décomposition additive unique (11) sera dit \emph{réductible} en les deux sous-groupes $\mA$ et $\mB$, ce que l'on écrit $\mG=\mA+\mB$; un groupe qui n'est pas réductible s'appelle \emph{irréductible}. Nous appelons aussi parfois les modules associés réductibles ou irréductibles. Il s'ensuit donc ce qui suit.

% unit: noether.p17.thm.001
\textbf{Théorème 1.}\footnote{Cf. Schmeidler, loc. cit., p. 39.} \emph{Un module $\mM$ est représentable comme plus petit multiple commun de deux diviseurs propres premiers entre eux $\mL$ et $\mN$ si et seulement si son groupe quotient $\mG$ est réductible en deux sous-groupes $\mA$ et $\mB$. Dans ce cas, $\mA$ et $\mB$ sont respectivement isomorphes aux groupes quotients de $\mL$ et de $\mN$.}
